\chapter{Securitate}

\section{Introducere}

Securitatea nu e o funcționalitate pe care o adaugi la final, ci ceva ce trebuie
avut în minte peste tot. De aceea, pentru fiecare formular, încărcare de fișier
și acțiune de administrare m-am întrebat același lucru: ce s-ar întâmpla dacă
cineva ar apela-o direct, fără să treacă prin interfață~\cite{owasp-top10}?

\section{Autentificare și parole}

Parolele nu sunt niciodată stocate în clar; gestionarea lor este delegată
componentei \texttt{PasswordHasher} din ASP.NET Identity~\cite{identity-docs},
care aplică \emph{hashing} cu \emph{salt} unic per utilizator, în linie cu
recomandările privind gestionarea autentificatorilor de tip memorat-secret
publicate de NIST~\cite{nist-800-63b}.

Autentificarea unui cont blocat este respinsă înainte ca sesiunea să fie creată.
Pentru a evita confirmarea existenței unui cont blocat unui atacator care nu
cunoaște parola corectă, verificarea parolei se face separat de autentificarea
efectivă -- mesajul de blocare apare exclusiv când parola introdusă este corectă:

\begin{lstlisting}[caption={Verificarea blocării unui cont, ulterior validării parolei}]
var passwordCheck = await _signInManager.CheckPasswordSignInAsync(
    user, Input.Password, lockoutOnFailure: false);

if (passwordCheck.Succeeded)
{
    if (user.BannedUntil.HasValue &&
        user.BannedUntil.Value > DateTime.UtcNow)
    {
        ModelState.AddModelError(string.Empty, banMessage);
        return Page();
    }
    await _signInManager.SignInAsync(user, isPersistent: false);
    return LocalRedirect(returnUrl);
}
\end{lstlisting}

\section{Autorizare bazată pe roluri}

Panoul de administrare, precum și fiecare acțiune sensibilă din
\texttt{AdminController} (blocare, ștergere, promovare la rolul de administrator,
gestionarea materiilor) sunt protejate prin \texttt{[Authorize(Roles = "Admin")]}
~\cite{identity-roles-docs}, verificat de \emph{middleware}-ul de autorizare al
ASP.NET Core, înaintea execuției oricărui cod din controller. Aceeași filosofie
a fost aplicată restricțiilor de domeniu care nu se mapează pe un rol Identity:
doar un utilizator cu \texttt{Role = "Teach"} poate atașa fișiere unei
conversații, verificare realizată în endpoint-ul de încărcare.

\section{Protecția formularului de înregistrare}

Formularul de înregistrare este protejat prin \emph{Google reCAPTCHA
v2}~\cite{recaptcha-docs}. Token-ul generat de widget-ul din pagină este trimis
către server și verificat printr-un apel către serviciul oficial de validare al
Google, înainte ca vreun cont să fie creat:

\begin{lstlisting}[caption={Verificarea token-ului reCAPTCHA (fragment)}]
public async Task<bool> VerifyAsync(
    string? recaptchaResponse, string? remoteIp)
{
    if (string.IsNullOrWhiteSpace(recaptchaResponse))
        return false;

    var response = await client.PostAsync(
        "https://www.google.com/recaptcha/api/siteverify",
        new FormUrlEncodedContent(parameters));

    var json = await response.Content.ReadAsStringAsync();
    using var doc = JsonDocument.Parse(json);
    return doc.RootElement.TryGetProperty("success", out var s)
           && s.GetBoolean();
}
\end{lstlisting}

Dacă cheia secretă nu este configurată, verificarea eșuează implicit (\emph{fail
closed}) -- înregistrarea rămâne blocată până la configurarea corectă, în loc să
permită tuturor înregistrărilor să treacă nevalidate.

\section{Validarea fișierelor încărcate}

Atât avatarele de profil, cât și fișierele atașate conversațiilor sunt validate
pe server, nu doar prin atributul \texttt{accept} al elementului HTML de tip
fișier (ușor de ocolit): tipul este verificat printr-o listă albă de extensii,
dimensiunea este limitată explicit, iar numele fișierului salvat pe disc este
generat aleator, pentru a preveni atât suprascrierea accidentală a unui fișier
existent, cât și deducerea, din nume, a provenienței acestuia.

\section{Restricționarea accesului la conversații}

O conversație există numai între doi utilizatori legați printr-o cerere de
colaborare acceptată. Interfața respectă regula în mod natural -- butonul de chat
apare exclusiv în lista de colaborări --, dar ea nu poate rămâne doar acolo: un
utilizator autentificat ar putea deschide manual adresa unei conversații cu orice
cont, indicând identificatorul acestuia în URL. Verificarea este, prin urmare,
repetată pe server:

\begin{lstlisting}[caption={Verificarea existenței unei colaborări acceptate}]
var areConnected = await _context.HireRequests.AnyAsync(h =>
    h.IsAccepted &&
    ((h.RequesterId == currentUserId && h.TargetId == userId) ||
     (h.RequesterId == userId && h.TargetId == currentUserId)));

if (!areConnected) return Forbid();
\end{lstlisting}

Aceeași condiție este aplicată atât în endpoint-ul care persistă un mesaj text,
cât și în cel de încărcare a fișierelor.

\section{Protecția împotriva atacurilor CSRF}

Un atac de tip \emph{Cross-Site Request Forgery} exploatează faptul că browser-ul
atașează automat \emph{cookie}-ul de sesiune oricărei cereri către aplicație,
inclusiv uneia inițiate de pe un site străin: o pagină ostilă poate trimite, în
numele unui administrator autentificat, o cerere de ștergere a unui
cont~\cite{owasp-top10}.

Contramăsura standard este un \emph{token} unic per sesiune, inclus în fiecare
formular și validat la primirea cererii. În ASP.NET Core, paginile Razor îl
validează automat, în timp ce controllerele MVC nu o fac decât dacă li se cere
explicit -- exact categoria de acțiuni sensibile din \texttt{AdminController} și
\texttt{HireController}. Validarea a fost, de aceea, activată global:

\begin{lstlisting}[caption={Activarea globală a validării token-ului antiforgery}]
builder.Services.AddAntiforgery(options =>
    options.HeaderName = "RequestVerificationToken");

builder.Services.AddControllersWithViews(options =>
{
    options.Filters.Add(new AutoValidateAntiforgeryTokenAttribute());
});
\end{lstlisting}

Formularele aplicației sunt generate prin \emph{tag helper}-ul de formular, care
include deja token-ul ca un câmp ascuns, deci nicio vizualizare nu a trebuit
modificată. Singurele două cereri care nu provin dintr-un formular -- trimiterea
unui mesaj și încărcarea unui fișier -- transmit token-ul printr-un antet HTTP.

\section{Prevenirea falsificării expeditorului unui mesaj}

Endpoint-ul API prin care se persistă un mesaj nu are voie să aibă încredere în
identificatorul de expeditor trimis de client; acesta este întotdeauna suprascris
cu identitatea utilizatorului autentificat:

\begin{lstlisting}[caption={Suprascrierea expeditorului cu identitatea autentificată}]
var senderId = _userManager.GetUserId(User);
if (senderId == null) return Unauthorized();

message.SenderId = senderId;
\end{lstlisting}

Fără această verificare, un utilizator autentificat ar fi putut construi manual
o cerere HTTP care să pară provenită de la un alt cont.

\section{Protecția integrității datelor și a rolului de administrator}

Ștergerea unui cont elimină explicit toate înregistrările care fac referire la
acel utilizator -- mesajele trimise și primite, cererile de colaborare și
codurile OTP -- înainte de ștergerea efectivă a contului, evitând înregistrările
orfane.

Un administrator nu se poate retrograda, bloca sau șterge pe sine, iar niciun
administrator nu poate fi blocat sau șters de un altul, pentru ca platforma să nu
rămână, accidental sau prin intervenție rău-intenționată, fără niciun cont capabil
să acceseze panoul.

\section{Confidențialitatea datelor utilizatorilor}

Platforma stochează date cu caracter personal (nume, email, biografie, avatar,
conținutul mesajelor), motiv pentru care principiile de minimizare prevăzute de
Regulamentul General privind Protecția Datelor~\cite{gdpr} au fost avute în vedere
la proiectare: utilizatorul își poate elimina avatarul oricând, iar ștergerea unui
cont înlătură și fișierul de pe disc, nu doar rândul din baza de date. Codurile OTP
expiră în 10 minute și sunt șterse automat, reducând fereastra în care ar putea fi
exploatate dacă baza de date ar fi compromisă.
